\documentclass[10pt]{article}
\usepackage[a4paper,margin=0.5in]{geometry}
\usepackage{amsmath,amssymb,amsfonts,mathtools}
\usepackage{xcolor}
\usepackage{array,longtable,booktabs}
\usepackage[colorlinks=true,linkcolor=blue]{hyperref}
\allowdisplaybreaks
\numberwithin{equation}{section}
\setlength{\jot}{2pt}
\setlength{\LTpre}{4pt}\setlength{\LTpost}{4pt}
\newcommand{\KK}{\mathcal K}
\newcommand{\NN}{\mathcal N}
\newcommand{\NO}{\mathcal N_0}
\newcommand{\LL}{L_\Lambda}
\newcommand{\Lk}{L_k}
\newcommand{\J}{\mathcal J}
\newcommand{\B}{\mathcal B}
\setcounter{tocdepth}{1}

\title{How the two-$\rho$ terms give JIMWLK and BFKL}
\date{}

\begin{document}
\maketitle
\tableofcontents

%==========================================================================================
\section{The result in short}\label{sec:short}
Notation, as in your notes:
\begin{equation}
\NO\equiv\frac{1}{(2\pi)^3}\,\frac{g^4}{16\pi^5}\,\frac1{k^+},\qquad \LL\equiv\log\frac\vee\Lambda,\qquad\Lk\equiv\log\frac\vee{k^+},\qquad \NN=\NO\LL .
\end{equation}
\begin{itemize}
\item $\J$ is JIMWLK$\otimes$LO, the JIMWLK evolution of the LO cross section per unit rapidity (Sec.~\ref{sec:J}).
\item $\B$ is BFKL$\otimes$LO, the BFKL evolution of the LO cross section per unit rapidity (Sec.~\ref{sec:B}).
\end{itemize}
The two-$\rho$ terms at leading logarithm are
\begin{equation}
\boxed{\;\frac{d^3N}{d^3k}\bigg|_{\rm LL}^{2\rho}=\NO\Big[\LL\;\J+\Lk\,\big(\B-\J\big)\Big]=\NO\Big[\log\frac{k^+}{\Lambda}\;\J+\log\frac{\vee}{k^+}\;\B\Big].\;}
\label{eq:result}
\end{equation}
This note shows the two halves explicitly.

\paragraph{1. JIMWLK: the coefficient of $\log(\vee/\Lambda)$} (Part~I). It is the sum of three sets of two-$\rho$ terms:
\begin{enumerate}
\item[(a)] the two-$\rho$ terms of the three-$\rho$ equation (\texttt{TwoRho\_from\_final.tex}, Eq.~(1), 335 terms);
\item[(b)] the genuine two-$\rho$ terms of \texttt{Evolution\_two\_rho.tex}, Eq.~\eqref{eq:gen} (@@NGEN@@ terms);
\item[(c)] the two-$\rho$ lists of \texttt{4rho.tex}, Eq.~\eqref{eq:four} (@@NFOUR@@ terms).
\end{enumerate}
Added up kernel structure by kernel structure, they give exactly JIMWLK$\otimes$LO, Eq.~\eqref{eq:Jexp}. This holds in all @@NSA@@ kernel structures and in all @@NBA@@ independent colour structures with a nonzero entry. The coefficients are exact polynomials in $N_c$.

\paragraph{2. BFKL: the coefficient of $\log\vee$} (Part~II). Both of your logarithms contain $\log\vee$. So the coefficient of $\log\vee$ is
\begin{equation*}
P=[\text{coefficient of }\log(\vee/\Lambda)]+[\text{coefficient of }\log(\vee/k^+)].
\end{equation*}
Row by row, $P$ is the integrand of the $p^+$ integral at its $\vee$ end, where the unobserved gluon is harder than the measured one. The $P$ two-$\rho$ terms of all rows are listed in Eq.~\eqref{eq:Plist} (@@NPL@@ terms). Their sum is exactly BFKL$\otimes$LO, Eq.~\eqref{eq:Bexp}, in all @@NSP@@ kernel structures and @@NBP@@ independent colour structures with a nonzero entry.

\paragraph{Why the split is like this} (shown in \texttt{TwoRho\_all.pdf}, Sec.~3, without taking any limit).
\begin{itemize}
\item $\Lambda$ enters a row only through the limits of its $p^+$ integral where a gluon carries momentum $\Lambda$. That gluon is softer than the measured one: it evolves the target (JIMWLK).
\item $\vee$ enters only through the upper limit, where the unobserved gluon is harder than the measured one: it evolves the projectile (BFKL).
\item Of your two logarithms only $\log(\vee/\Lambda)$ contains $\Lambda$. So the coefficient of $\log(\vee/\Lambda)$ is JIMWLK. BFKL is the coefficient of $\log\vee$, which sits in both of your logarithms.
\end{itemize}

\paragraph{Variables.}
\begin{itemize}
\item The measured gluon is at $w$ (amplitude) and $w'$ (c.c.); the unobserved gluon is at $z$; the two charges are at $x$ and $x'$.
\item $x$ is tied to $w$ and $x'$ to $w'$ by their emission kernels. If one charge emits at both $w$ and $w'$, it is $x'$.
\item $\KK(X)=X/X^2$. The phase $e^{-ik(w'-w)}$ and the integrals over $x,x',z$ are understood.
\item The two charges are a symmetric product, $\rho^a(p)\rho^b(q)\equiv\frac12\{\rho^a(p),\rho^b(q)\}$.
\item The colour algebra uses $U^{bc}=U^{cb}$.
\end{itemize}

%==========================================================================================
\part{JIMWLK: the coefficient of \texorpdfstring{$\log(\vee/\Lambda)$}{log(vee/Lambda)}}
\section{The three sets of two-\texorpdfstring{$\rho$}{rho} terms}\label{sec:sets}
\subsection{(a) From the three-\texorpdfstring{$\rho$}{rho} equation}
These are the 335 terms of \texttt{TwoRho\_from\_final.tex}, Eq.~(1). They are obtained by reordering every black and blue line of \texttt{ThreeRho\_final.tex}. They are used here exactly as written there.

\subsection{(b) The genuine two-\texorpdfstring{$\rho$}{rho} terms}
These are your terms of \texttt{Evolution\_two\_rho.tex}, in the same notation.
\begin{itemize}
\item \emph{Group I, Row II.} The bracket $-\frac2\epsilon-2\gamma-\log\frac{(x-w)^2\mu^2}{4}$ is $\frac1\pi\int_z[K(x,w,z)-K(w,w,z)]$ (fix F6). Only its pole and constant change: $\frac1\pi\int_zK(a,b,z)=\frac1{\bar\epsilon}-\log(4\pi^2\mu^2(a-b)^2)$ in dimensional regularization. We use the integrand.
\item \emph{Group I, Row V.} Parts 2a and 2b together. The $\log(\vee/\Lambda)$ part of 2b has the same two kernels as 2a, and the $N_cU^{af}(x)$ terms cancel.
\item \emph{Group II, Row IV.} The bracket $\frac1{\bar\epsilon}-\log(4\pi^2\mu^2(w-w')^2)$ is written as the integrand $\frac1\pi[K(w,w',z)-\frac12K(w,w,z)-\frac12K(w',w',z)]$, which has the same integral.
\item \emph{Group III, Row III, part 2.} The overall sign is reversed (fix F5).
\end{itemize}
After the colour algebra, and collecting equal terms within each row:
@@GENA@@

\subsection{(c) The two-\texorpdfstring{$\rho$}{rho} lists of \texttt{4rho.tex}}
These are your lists ``Two $\rho$ contribution'' of \texttt{4rho.tex}, with fixes F1 (drop the $-8N_c$ term) and F2 (the repeated index in $dN^{(1)}$, term 2). Collected:
@@FOURA@@
The complete two-$\rho$ content of the four-$\rho$ terms is this list plus the two-$\rho$ terms of the remainders $d^3N_a$, $d^3N_b$. The remainders are part of set (a).

%==========================================================================================
\section{JIMWLK\texorpdfstring{$\otimes$}{x}LO}\label{sec:J}
\paragraph{Operator form.} With $\delta Y=\LL$, the JIMWLK evolution of the LO cross section is
\begin{align}
\NN\,\J={}&-\NN\int e^{-ik\cdot(w'-w)}\int_{x,x'}\KK(x'-w')\cdot\KK(x-w)\;\tilde H\Big\{\rho^{b'}(x')\big[U(x')-U(w')\big]^{ab'}\big[U(x)-U(w)\big]^{ab}\rho^b(x)\Big\},\notag\\
\tilde H={}&\int_{u,v,z}K(u,v,z)\Big\{J_R^h(u)J_R^h(v)+J_L^m(u)J_L^m(v)-U^{mh}(z)\big[J_R^h(u)J_L^m(v)+J_L^m(u)J_R^h(v)\big]\Big\},
\label{eq:Jop}
\end{align}
with $K(u,v,z)=\KK(u-z)\cdot\KK(v-z)$ and the definitions below.
\begin{itemize}
\item $J_R^h(u)\,U(p)=\delta^{(2)}(u-p)\,U(p)T^h$, $J_L^m(u)=U^{mg}(u)J_R^g(u)$, $(T^h)_{bc}=-if^{hbc}$.
\item This is $H_{\rm JIMWLK}$ of your Eq.~(HLO), with $d/dY=-H$.
\item $\tilde H$ differentiates the Wilson lines only. So $u$ and $v$ are each one of the points $x,\ x',\ w,\ w'$ where the LO cross section has a Wilson line: the soft gluon at $z$ connects two of them.
\item This gives the kernel structures $K(x,x,z)$, $K(x',x',z)$, $K(x,x',z)$ (soft gluon between the charges), $K(w,w,z)$, $K(w',w',z)$, $K(w,w',z)$ (between the measured gluons), and the mixed $K(x,w',z)$, $K(x',w,z)$.
\end{itemize}
\paragraph{Written out.} Carrying out the derivatives and the colour algebra:
@@JEQ@@
The terms of the structures that do not appear here cancel within the three sets.

%==========================================================================================
\section{The sum, structure by structure}\label{sec:Jsum}
\subsection{How the comparison is done}
\begin{enumerate}
\item Every term of (a), (b), (c) is written with its kernel in one of the @@NSA@@ structures $S_n$ of Table~\ref{tab:SA}.
\item Within a structure, the colour structures are not all independent: Jacobi-type identities relate them (App.~\ref{app:rel}). Each colour structure is written in a basis of independent ones, with coefficients that are exact polynomials in $N_c$ (fixed from SU(3) and SU(4)).
\item For each basis colour structure, the coefficients from (a), (b), (c) are added and compared with $\J$.
\end{enumerate}
\begin{table}[h]
\centering\small
\begin{tabular}{@{}l l c@{}}
& kernel structure & in $\J$\\ \hline
@@STRUCTA@@
\end{tabular}
\caption{The kernel structures of the $\log(\vee/\Lambda)$ two-$\rho$ terms.}
\label{tab:SA}
\end{table}

\subsection{Worked example 1: the soft gluon between the two charges}
@@WA@@

\subsection{Worked example 2: a structure that is not in JIMWLK}
@@WA2@@

\subsection{All structures}
Columns: the coefficient from (a) the three-$\rho$ equation, (b) the genuine terms, (c) \texttt{4rho.tex}, their sum, and the coefficient in $\J$. All are in units of $\NN$, multiplying [kernel]$\times$[colour structure].
@@JTAB@@

%==========================================================================================
\part{BFKL: the coefficient of \texorpdfstring{$\log\vee$}{log vee}}
\section{What has to be added}\label{sec:P}
Write every row as
\begin{equation*}
\frac{d^3N_{\rm row}}{d^3k}=\NO\big[A_{\rm row}\,\LL+B_{\rm row}\,\Lk\big]+(\text{no large log}),\qquad P_{\rm row}=A_{\rm row}+B_{\rm row}.
\end{equation*}
\begin{itemize}
\item $A_{\rm row}$ are the $\log(\vee/\Lambda)$ terms of Part~I.
\item $B_{\rm row}$ are the $\log(\vee/k^+)$ terms of the same row. In your notes they come from functions such as $\log(\Lambda/k^+)$, $\log\frac\vee{\Lambda+k^+}$ and $\log(\frac\vee{k^+}-1)$, or hidden in the plus-prescription integrals.
\item $P_{\rm row}$ is the coefficient of $\log\vee$: the integrand of the row at the $\vee$ end of its $p^+$ integral.
\end{itemize}
For the two-$\rho$ terms:
\begin{itemize}
\item The two-$\rho$ terms from the four-$\rho$ terms (the remainders $d^3N_a$, $d^3N_b$ and \texttt{4rho.tex}) have $P=A$: their $p^+$ integral is $\int_\Lambda^\vee dp^+/p^+$ with a $p^+$-independent factor.
\item Group I Row V, $\mathcal D$ and the genuine Group I terms have $P=0$: their $p^+$ integral runs from $\Lambda$ to $k^+-\Lambda$, so it has no $\vee$.
\item Group I Row II has only the $\mathfrak F_\alpha$ part, with the sign reversed (fix F4: this part multiplies $-\log(\vee/k^+)$).
\item Every other row has its $\vee$-end integrand. Fix P1 (Group I Row III) and fix F5 (Group III Row III, part 2) apply.
\end{itemize}
The symmetrized three-$\rho$ parts of $P$ cancel (the $\log(\vee/k^+)$ note), so only two-$\rho$ terms remain:
@@PLIST@@
The genuine Group II Row IV term is the $\vee$ end of the same $p^+$ integral whose $\Lambda$ end gives the genuine term of Part~I:
\begin{itemize}
\item at the $\Lambda$ end, the soft gluon at $z$ is emitted by the measured gluon ($K(w,w',z)$, JIMWLK);
\item at the $\vee$ end, the gluon at $z$ is the parent and emits the measured gluon at $w,w'$ ($\KK(w-z)\cdot\KK(w'-z)\,K(x,x',z)$, BFKL).
\end{itemize}

%==========================================================================================
\section{BFKL\texorpdfstring{$\otimes$}{x}LO}\label{sec:B}
\paragraph{Operator form.} Let $O(x',x)=\rho^{b'}(x')\,M^{b'b}(x',x)\,\rho^b(x)$ with $M^{b'b}(p',p)=\big[U(w')-U(p')\big]^{ab'}\big[U(w)-U(p)\big]^{ab}$. The LO cross section is $\propto\int\KK(x'-w')\cdot\KK(x-w)\langle O(x',x)\rangle$. A gluon harder than the measured one is added to the projectile: its charge $j^a(z)$ adds to $\rho^a$, and the valence charges get the virtual correction. Then
\begin{align}
\B={}&2\int_{x,x',u,v,z}e^{-ik\cdot(w'-w)}\,K(u,v,z)\,\Big\{\notag\\
&-\tfrac12\,\KK(x'-w')\cdot\KK(x-w)\,\Big[\rho^h(u)\rho^h(v)\,O(x',x)-2\rho^h(u)\,O(x',x)\,\rho^h(v)+O(x',x)\,\rho^h(u)\rho^h(v)\Big]\notag\\
&+\KK(z-w')\cdot\KK(z-w)\;\rho^{e'}(u)\rho^e(v)\,(T^{b'}T^b)_{e'e}\,M^{b'b}(z,z)\notag\\
&+\KK(x'-w')\cdot\KK(z-w)\;\rho^{e'}(u)\rho^{b'}(x')\rho^e(v)\,(T^b)_{e'e}\,M^{b'b}(x',z)\notag\\
&+\KK(z-w')\cdot\KK(x-w)\;\rho^{e'}(u)\rho^{b}(x)\rho^e(v)\,(T^{b'})_{e'e}\,M^{b'b}(z,x)\Big\}.
\label{eq:Bop}
\end{align}
\begin{itemize}
\item \emph{First line:} the virtual correction of the valence charges (dipole kernel). Here $u$, $v$ are charges.
\item \emph{Second line:} both measured gluons are emitted by the gluon at $z$.
\item \emph{Last two lines:} one measured gluon is emitted by a valence charge and the other by the gluon at $z$.
\end{itemize}
In the colour-singlet projection this is Eq.~(LOBFKL) of your Appendix E. After the charges are reordered, its three- and four-$\rho$ parts vanish. Written out:
@@BEQ@@

%==========================================================================================
\section{The sum, structure by structure}\label{sec:Bsum}
Same method as Sec.~\ref{sec:Jsum}, now for $P$ against $\B$. The columns are the $\vee$-end terms of the three-$\rho$ rows, of the genuine rows, and of \texttt{4rho.tex} ($P=A$).
\begin{table}[h]
\centering\small
\begin{tabular}{@{}l l c@{}}
& kernel structure & in $\B$\\ \hline
@@STRUCTP@@
\end{tabular}
\caption{The kernel structures of the $\log\vee$ two-$\rho$ terms.}
\label{tab:SP}
\end{table}

\subsection{Worked example 1: both measured gluons from the gluon at \texorpdfstring{$z$}{z}}
@@WP@@
This whole structure comes from the $\vee$ end of the genuine Group II, Row IV term (Sec.~\ref{sec:P}). It is the second line of Eq.~\eqref{eq:Bop}. It has no counterpart in Part~I: no $\log(\vee/\Lambda)$ term has the measured gluon emitted by the gluon at $z$.

\subsection{Worked example 2: the soft gluon between the two charges}
@@WP2@@

\subsection{All structures}
@@BTAB@@

%==========================================================================================
\section{Putting the two together}\label{sec:together}
\begin{itemize}
\item Part~I: $\sum_{\rm rows}A_{\rm row}=\J$, the coefficient of $\LL=\log(\vee/\Lambda)$.
\item Part~II: $\sum_{\rm rows}P_{\rm row}=\B$, the coefficient of $\log\vee$.
\item Since $P=A+B$, the coefficient of $\Lk=\log(\vee/k^+)$ is $\sum_{\rm rows}B_{\rm row}=\B-\J$.
\end{itemize}
This is Eq.~\eqref{eq:result}:
\begin{equation*}
\frac{d^3N}{d^3k}\bigg|_{\rm LL}^{2\rho}=\NO\Big[\log\frac{\vee}{\Lambda}\;\J+\log\frac{\vee}{k^+}\,\big(\B-\J\big)\Big]=\NO\Big[\log\frac{k^+}{\Lambda}\;\J+\log\frac{\vee}{k^+}\;\B\Big].
\end{equation*}
The target evolves (JIMWLK) over the rapidity interval below the measured gluon, $\log(k^+/\Lambda)$. The projectile evolves (BFKL) over the interval above it, $\log(\vee/k^+)$. Together these make $\delta Y=\log(\vee/\Lambda)$.

%==========================================================================================
\appendix
\section{Colour identities used in the basis reduction}\label{app:rel}
In each kernel structure where the colour structures are linearly dependent, every dependent colour structure equals the following combination of basis structures. The identities are exact and were checked for SU(3) and SU(4).
\subsection{Part I}
@@RELA@@
\subsection{Part II}
@@RELP@@

\section{Scripts}
\begin{itemize}
\item \texttt{evo\_build.py}: the three sets and the targets, by kernel structure and colour structure, and the basis reduction.
\item \texttt{evo\_note.py}: writes this file from \texttt{evo\_static.tex}.
\end{itemize}
Set (a) is read from \texttt{ThreeRho\_final.tex} (black and blue lines). Sets (b), (c), the $P$ terms and the targets come from the scripts of \texttt{../two\_rho\_all}, \texttt{../two\_rho} and \texttt{../bfkl}.
\end{document}
